\documentclass{article}
\usepackage[T1]{fontenc}
\usepackage{lmodern}
\usepackage{graphicx}
\usepackage{amsmath,amssymb}
\usepackage[a4paper,margin=2cm]{geometry}
\usepackage[absolute,overlay]{textpos}
\setlength{\TPHorizModule}{1mm}
\setlength{\TPVertModule}{1mm}
\usepackage[indonesian]{babel}
\usepackage{hyperref}
\setlength{\emergencystretch}{2em}
% unit: Halaman 37

\begin{document}

% === Halaman 37 (python/native) ===

Elliptic Curve Cryptography (ECC) *)

• ECC adalah sistem kriptografi kunci-publik, sejenis dengan RSA, Rabin, 

ElGamal, D-H, dll.

• Setiap pengguna memiliki kunci publik dan kunci privat

     -  Kunci publik untuk enkripsi atau untuk verifikasi tanda  tangan digital

     -  Kunci privat untuk dekripsi atau untuk menghasilkan tanda tangan digital

• Kurva eliptik digunakan sebagai perluasan sistem kriptografi kunci-publik 

yang  lain:


1. Elliptic Curve Elgamal (ECEG)


2. Elliptic Curve Digital Signature (ECDSA)


3. Elliptic Curve Diffie-Hellman (ECDH)

37

*) Sumber bahan: Debdeep Mukhopadhyay, Elliptic Curve Cryptography ,

 Dept of Computer Sc and Engg IIT Madras

\end{document}
